\documentclass[12pt, a4paper]{article}

% ── Paquetes ──────────────────────────────────────────────────────────────────
\usepackage[utf8]{inputenc}
\usepackage[T1]{fontenc}
\usepackage[spanish]{babel}
\usepackage{amsmath, amssymb}
\usepackage{geometry}
\usepackage{hyperref}
\usepackage{enumitem}
\usepackage{titlesec}
\usepackage{lmodern}
\usepackage{booktabs}
\usepackage{float}

\geometry{margin=2.5cm}

\hypersetup{colorlinks=true, linkcolor=blue, urlcolor=blue, citecolor=blue}

% ── Portada ───────────────────────────────────────────────────────────────────
\title{%
  \LARGE \textbf{Series de Tiempo} \\[0.4em]
  \large Fundamentos y Aplicaciones \\[0.4em]
  \normalsize Repositorio Académico
}
\author{}
\date{}

% ══════════════════════════════════════════════════════════════════════════════
\begin{document}

\maketitle
\tableofcontents
\newpage

% ══════════════════════════════════════════════════════════════════════════════
\section{Descripción General}

Repositorio académico para el estudio teórico y aplicado de series de tiempo,
con énfasis en modelos clásicos y financieros. El objetivo es comprender los
procesos estocásticos, la identificación, estimación y validación de modelos,
y su aplicación en datos económicos y financieros.

% ══════════════════════════════════════════════════════════════════════════════
\section{Contenido del Curso}

\subsection{Introducción a Series de Tiempo}

\begin{itemize}
  \item ¿Qué es una serie de tiempo?
  \item Componentes:
    \begin{itemize}
      \item Tendencia
      \item Estacionalidad
      \item Ciclo
      \item Ruido
    \end{itemize}
  \item Estacionariedad:
    \begin{itemize}
      \item Estricta
      \item Débil
    \end{itemize}
  \item Transformaciones comunes:
    \begin{itemize}
      \item Diferenciación
      \item Logaritmos
      \item Desestacionalización
    \end{itemize}
\end{itemize}

% ══════════════════════════════════════════════════════════════════════════════
\subsection{Procesos Estocásticos y Modelos ARIMA}

\subsubsection{Modelos ARIMA$(p,d,q)$}

\begin{itemize}
  \item \textbf{AR (Autorregresivo)}
  \item \textbf{MA (Media Móvil)}
  \item \textbf{Integración (diferenciación)}
  \item Interpretación económica y estadística
  \item Condiciones de estacionariedad e invertibilidad
\end{itemize}

\subsubsection{Modelos SARIMA$(p,d,q)(P,D,Q)_s$}

\begin{itemize}
  \item Extensión estacional de ARIMA
  \item Componentes estacionales
  \item Periodicidad $s$
  \item Casos típicos (mensual, trimestral)
\end{itemize}

% ══════════════════════════════════════════════════════════════════════════════
\subsection{Identificación y Diagnóstico}

\begin{itemize}
  \item Función de autocorrelación (ACF)
  \item Función de autocorrelación parcial (PACF)
  \item Correlogramas
  \item Diagnóstico de residuos
\end{itemize}

% ══════════════════════════════════════════════════════════════════════════════
\subsection{Estimación y Selección de Modelos}

\begin{itemize}
  \item Función de log-verosimilitud
  \item ECF (Empirical Characteristic Function)
  \item Criterios de información:
    \begin{itemize}
      \item \textbf{AIC} (Akaike Information Criterion)
      \item \textbf{BIC} (Bayesian Information Criterion)
    \end{itemize}
  \item Comparación de modelos
\end{itemize}

% ══════════════════════════════════════════════════════════════════════════════
\subsection{Series de Tiempo Financieras}

\subsubsection{Características de series financieras}

\begin{itemize}
  \item Rendimientos
  \item Volatilidad
  \item Clustering de volatilidad
  \item No normalidad
  \item Leptocurtosis
\end{itemize}

% ══════════════════════════════════════════════════════════════════════════════
\subsection{Modelos de Volatilidad}

\subsubsection{Modelo GARCH}

\begin{itemize}
  \item ARCH
  \item GARCH$(p,q)$
  \item Interpretación financiera
  \item Persistencia de volatilidad
\end{itemize}

\subsubsection{Modelo EGARCH}

\begin{itemize}
  \item Efecto apalancamiento
  \item Asimetría
  \item Ventajas frente al GARCH clásico
\end{itemize}

% ══════════════════════════════════════════════════════════════════════════════
\section{Objetivos del Repositorio}

\begin{enumerate}
  \item Entender la teoría detrás de los modelos de series de tiempo.
  \item Aplicar modelos ARIMA/SARIMA a datos reales.
  \item Modelar y analizar volatilidad financiera.
  \item Comparar modelos usando criterios estadísticos.
  \item Desarrollar intuición econométrica.
\end{enumerate}

% ══════════════════════════════════════════════════════════════════════════════
\section{Tecnologías y Herramientas}

\begin{itemize}
  \item Python / R
  \item NumPy / pandas
  \item \texttt{statsmodels}
  \item \texttt{arch}
  \item Jupyter Notebooks
\end{itemize}

% ══════════════════════════════════════════════════════════════════════════════
\section{Estructura Sugerida del Repositorio}

\begin{table}[H]
  \centering
  \caption{Organización del repositorio}
  \label{tab:estructura}
  \begin{tabular}{ll}
    \toprule
    \textbf{Carpeta / Archivo} & \textbf{Contenido} \\
    \midrule
    \texttt{notebooks/}        & Jupyter Notebooks de cada tema \\
    \texttt{data/}             & Conjuntos de datos utilizados \\
    \texttt{scripts/}          & Funciones y utilidades reutilizables \\
    \texttt{reports/}          & Reportes y documentación en PDF/LaTeX \\
    \texttt{README.md}         & Este documento \\
    \bottomrule
  \end{tabular}
\end{table}

\end{document}
